\section{版本补充：模式匹配与联合类型}\label{rev:python:language}
本节的模式匹配与竖线联合类型从 Python 3.10 起可用；后面的异常组、TaskGroup 和 timeout 示例要求 Python 3.11 或更新版本。随书示例的本次运行环境为 Python 3.12。读者可先用 \code{python --version} 核对版本。

\code{match} 按顺序尝试模式，可同时检查结构并绑定字段。字典模式默认允许额外键；普通名字通常是捕获变量，不是与同名常量比较。条件守卫用 \code{if} 追加判断。由于 \code{bool} 是 \code{int} 的子类，示例先匹配布尔值，避免把 True 当作整数一。
\lstinputlisting[title={\code{examples/revision/modern_match.py}}]{examples/revision/modern_match.py}
预期输出依次为 \code{nonnegative: 7}、\code{note: Calculus}、\code{boolean}。形如 \code{int | str} 的联合类型表达“允许多种类型”，并不会自动转换输入，也不会代替运行时校验。本例把字典也列入参数类型，是为了与实际接收的数据结构一致；更细的工程接口可用 TypedDict 表达字段约定。

\subsection{异常组与 except*：按异常类型分派}
Python 3.11 引入 ExceptionGroup，用于把多项独立失败一起报告。\code{except*} 对匹配的子组处理，多个分支可在同一次异常处理中执行；未匹配部分仍会传播，不能认为处理一个子组就处理了全部失败。同一 try 不混写普通 except 与 except*，except* 子句中也不能用 return、break 或 continue。
\lstinputlisting[title={\code{examples/revision/exception_groups.py}}]{examples/revision/exception_groups.py}
脚本显式捕获两类问题，输出两行计数一。它是受控的异常组演示，与附录中故意无法正确执行的错误脚本分开。
语言细节可查 \href{https://docs.python.org/3/reference/compound_stmts.html}{Python 复合语句官方文档}。
